\hyperref[fig:5.8]{图 5.8} 进行比较. 添加一个特征后给出的 $T^{2}$ 控制图是否改变了你对过程稳定性的结论? 请解释.

\medskip
\noindent\mvbold{5.27}\quad 用\hyperref[tab:5.8]{表 5.8} 中 $x_{3}$ 与 $x_{4}$ 的数据, 构造未来观测值 $\mathbf{x}'=(x_{3},x_{4})$ 的预测椭圆. 注意, 预测椭圆必须从稳定过程出发求得. 对所得结果作出解释.

\medskip
\noindent\mvbold{5.28}\quad 作为金属片装配过程研究的一部分, 一个大的汽车制造商在汽车的六个位置使用传感器来记录与名义上的厚度的偏离 (mm). 前四个是当汽车主体完成时测量的, 后两个是装配早期在汽车底部测量的. 基于 50 辆汽车的数据见\hyperref[tab:5.14]{表 5.14}.

\noindent (a) 装配过程对于前 30 个看起来是稳定的. 用这些数据来估计 $\mathbf{S}$ 和 $\bar{\mathbf{x}}$. 然后用所有的数据 (50 个) 构造一个 $T^{2}$ 控制图.

\noindent (b) 哪一个位置看起来是需要找原因的?

\begin{mvtable}
\phantomsection\label{tab:5.14}\mvtabcap{表 5.14\quad 车体装配数据}
\centering
{\small
\begin{tabular}{|r|r|r|r|r|r|r|}
\hline
编号 & $x_{1}$ & $x_{2}$ & $x_{3}$ & $x_{4}$ & $x_{5}$ & $x_{6}$\\
\hline
$1$ & $-0.12$ & $0.36$ & $0.40$ & $0.25$ & $1.37$ & $-0.13$\\
$2$ & $-0.60$ & $-0.35$ & $0.04$ & $-0.28$ & $-0.25$ & $-0.15$\\
$3$ & $-0.13$ & $0.05$ & $0.84$ & $0.61$ & $1.45$ & $0.25$\\
$4$ & $-0.46$ & $-0.37$ & $0.30$ & $0.00$ & $-0.12$ & $-0.25$\\
$5$ & $-0.46$ & $-0.24$ & $0.37$ & $0.13$ & $0.78$ & $-0.15$\\
$6$ & $-0.46$ & $-0.16$ & $0.07$ & $0.10$ & $1.15$ & $-0.18$\\
$7$ & $-0.46$ & $-0.24$ & $0.13$ & $0.02$ & $0.26$ & $-0.20$\\
$8$ & $-0.13$ & $0.05$ & $-0.01$ & $0.09$ & $-0.15$ & $-0.18$\\
$9$ & $-0.31$ & $-0.16$ & $-0.20$ & $0.23$ & $0.65$ & $0.15$\\
$10$ & $-0.37$ & $-0.24$ & $0.37$ & $0.21$ & $1.15$ & $0.05$\\
$11$ & $-1.08$ & $-0.83$ & $-0.81$ & $0.05$ & $0.21$ & $0.00$\\
$12$ & $-0.42$ & $-0.30$ & $0.37$ & $-0.58$ & $0.00$ & $-0.45$\\
$13$ & $-0.31$ & $0.10$ & $-0.24$ & $0.24$ & $0.65$ & $0.35$\\
$14$ & $-0.14$ & $0.06$ & $0.18$ & $-0.50$ & $1.25$ & $0.05$\\
$15$ & $-0.61$ & $-0.35$ & $-0.24$ & $0.75$ & $0.15$ & $-0.20$\\
$16$ & $-0.61$ & $-0.30$ & $-0.20$ & $-0.21$ & $-0.50$ & $-0.25$\\
$17$ & $-0.84$ & $-0.35$ & $-0.14$ & $-0.22$ & $1.65$ & $-0.05$\\
$18$ & $-0.96$ & $-0.85$ & $0.19$ & $-0.18$ & $1.00$ & $-0.08$\\
$19$ & $-0.90$ & $-0.34$ & $-0.78$ & $-0.15$ & $0.25$ & $0.25$\\
$20$ & $-0.46$ & $0.36$ & $0.24$ & $-0.58$ & $0.15$ & $0.25$\\
$21$ & $-0.90$ & $-0.59$ & $0.13$ & $-0.58$ & $0.60$ & $-0.08$\\
$22$ & $-0.61$ & $-0.50$ & $-0.34$ & $-0.58$ & $0.95$ & $-0.08$\\
$23$ & $-0.61$ & $-0.20$ & $-0.58$ & $-0.20$ & $1.10$ & $0.00$\\
$24$ & $-0.46$ & $-0.30$ & $-0.10$ & $-0.10$ & $0.75$ & $-0.10$\\
$25$ & $-0.60$ & $-0.35$ & $-0.45$ & $0.37$ & $1.18$ & $-0.30$\\
$26$ & $-0.60$ & $-0.36$ & $-0.34$ & $-0.11$ & $1.68$ & $-0.32$\\
$27$ & $-0.31$ & $0.35$ & $-0.45$ & $-0.10$ & $1.00$ & $-0.25$\\
$28$ & $-0.60$ & $-0.25$ & $-0.42$ & $0.28$ & $0.75$ & $0.10$\\
$29$ & $-0.31$ & $0.25$ & $-0.34$ & $-0.24$ & $0.65$ & $0.10$\\
\hline
\end{tabular}}
\end{mvtable}
